\section{Conclusion}
This review has examined the design space of digital correlators and real-world implementations, leading to a clear rationale for a Python-based FX correlator in a four-element radio interferometer. In correlator design, two canonical architectures predominate. An \emph{XF correlator} first cross-multiplies time-domain signals and then applies a Fourier transform, whereas an \emph{FX correlator} first Fourier-transforms each antenna signal and then computes cross-products in the frequency domain. Although the two methods yield mathematically equivalent spectral outputs, they differ in computational structure. In practice, the FX approach takes advantage of fast Fourier transforms (FFT) with cost \( \mathcal{O}(N \log N) \) per channel, making it more efficient for wideband signals, while XF relies on direct convolutions in time (cost \( \mathcal{O}(N^2) \) for \( N \) lags)~\cite{Iguchi2008}. To further balance these trade-offs, hybrid schemes have been proposed: for example, the FFX correlator performs a coarse FFT filterbank followed by fine-grained FX processing~\cite{Iguchi2008}. Such two-stage architectures (also called FXF or polyphase FX) can reduce data rates early on while preserving frequency resolution. In practice, the choice of architecture is driven by factors such as antenna count, bandwidth, channelisation needs, and hardware capabilities.

Real-world observatories illustrate these choices. ALMA (Atacama Large Millimeter Array) initially planned a pure XF design but adopted an FXF correlator~\cite{Escoffier2007}. ALMA’s implementation uses tunable digital anti-aliasing filters for coarse sub-band splitting and then an FX correlator for fine spectral channels~\cite{Escoffier2007}. For large \( N \) and multi-GHz bandwidth, this hybrid FXF design meets ALMA’s requirements of high spectral resolution and massive data throughput. In contrast, the MeerKAT telescope (64 antennas) employs a custom FPGA-based FX correlator (the SKARAB boards) for its 856 MHz instantaneous bandwidth~\cite{Callanan2020}. Each SKARAB board hosts four X-engines, ingesting 27.2 Gbps via 40~GbE links, with channelisation performed in hardware. Callanan~\cite{Callanan2020} demonstrated that modern GPUs can substitute for these SKARAB X-engines. In a prototype, a single 1080Ti GPU (using the xGPU library) processed two SKARABs’ worth of data, at a fraction of the cost (R143k vs R490k). The GPU-based design occupied half the rack space, and although it doubled the power draw (400W vs 180W), it delivered comparable correlation results. This MeerKAT study highlights that GPU-FPGA hybrids can yield flexible correlators with favourable cost and density trade-offs~\cite{Callanan2020}.

The DiFX project illustrates software FX correlators in practice. Designed for VLBI, DiFX replaced decades of ASIC correlators to provide a highly flexible software FX correlator~\cite{Deller2011}. Deller et al.~\cite{Deller2011} report that DiFX supports arbitrary antenna configurations, bandwidths, and experiment modes in VLBI, exploiting clusters of general-purpose CPUs. For example, the VLBA correlator runs DiFX on a 68-node cluster (1360 CPU cores) for ten-station observing, taking advantage of commodity hardware. DiFX’s success demonstrates that for small to moderate antenna counts (\( N \sim 10 \)) and limited real-time constraints, a software FX correlator can be cost-effective and extensible. Similarly, the CHIME pathfinder telescope (256 inputs) adopted a mixed FX scheme: F-engine channelisation on FPGAs and X-engine cross-multiplication on GPUs~\cite{Denman2015}. Denman et al.~\cite{Denman2015} describe CHIME’s X-engine built from consumer-grade AMD GPUs using OpenCL. This system achieved 10.5~TOPS of processing for 400~MHz of bandwidth (32768 baselines), illustrating that commodity GPUs can handle large-scale correlation for fixed-baseline arrays. In all these cases, FX or hybrid correlators dominate, leveraging modern FFT techniques and parallel hardware. The architecture comparisons in the literature thus show that FX architectures generally offer superior scalability with bandwidth (via the FFT cost of \( \mathcal{O}(N\log N) \)) and easier channel isolation, whereas XF is occasionally preferred only for very small-channel or legacy systems~\cite{Escoffier2007, Denman2015, Iguchi2008}.

In light of these findings, an FX correlator implemented entirely in software is justified for the target four-element array. With only four antennas (six baselines), the total computational demand is modest: for example, the number of baselines is \( N(N-1)/2 = 6 \). Even at high sample rates, the data volume remains relatively low. For reference, Gorthi~\cite{Gorthi2021} achieved an aggregated input bandwidth of \(\sim450\)~Mbps in a 12-antenna testbed using FPGA digitizers. Scaling to four antennas suggests data rates on the order of a few hundred Mbps, which is easily carried by standard 1~Gbps or 10~Gbps networks with conservative buffering. Thus, bandwidth and I/O constraints are readily managed with modest hardware. Because the four-element system has no extreme throughput demands, a software solution need not be as heavily optimised as on a VLBI cluster or SKA-scale array. The literature notes that even for real-time pipelines (e.g.\ FRB searches), correlator data flows can be sustained by large ring buffers or memory-mapped I/O to absorb network jitter~\cite{Leung2024, Merry2023}.

The selection of Python as the implementation language flows naturally from these constraints. Python has become ubiquitous in astronomical software due to its rapid development capability and rich ecosystem~\cite{Gorthi2021}. Specifically, Python enables concise, expressive code for high-level tasks (data acquisition, control flow, I/O, and scheduling) while deferring heavy lifting to optimised libraries. The reviewed literature highlights several Python advantages: its ecosystem includes scientific libraries such as NumPy, SciPy, and Astropy, as well as GPU interfaces like CuPy and PyCUDA~\cite{Merry2023, Gorthi2021}. These tools allow a correlator developer to leverage well-tested algorithms (FFTs, complex linear algebra, HDF5 I/O) directly. Moreover, Python integrates seamlessly with C/C++/CUDA code via ctypes, Cython, or Python/C APIs. For example, existing projects use Python to control FFT routines (via CuPy or FFTW bindings) and to launch GPU kernels for cross-multiplication~\cite{Romein2021, Callanan2020}.

We must also consider performance optimisation. Python itself cannot perform millions of multiply-accumulate operations per second in pure interpreted code; as the literature emphasises, performance-critical loops are offloaded. In a Python FX correlator, the F-engine (channeliser) can use GPU-accelerated FFTs (CuPy, PyCUDA, or even FFTW via Cython), while the X-engine (baselines) can call a GPU-accelerated cross-multiply library (e.g.\ xGPU~\cite{Callanan2020} or custom CUDA code). The combination of batch processing and GPU parallelism has proven highly effective. For instance, PyFX~\cite{Leung2024} employs NumPy and Numba to handle FFTs and correlation on GPUs, meeting CHIME/FRB’s real-time constraints. Similarly, Merry~\cite{Merry2023} implemented a real-time channeliser on an NVIDIA V100 using CUDA and Python for control, achieving 2 GSps × 4 streams throughput. These examples show that Python-driven correlators can achieve performance comparable to lower-level systems, provided optimised libraries and large data blocks are used. Additionally, new hardware features further boost performance: Romein~\cite{Romein2021} introduces a tensor-core correlator library that exploits NVIDIA Tensor Cores for mixed-precision X-engines, yielding a 5–10× speedup over standard GPU cores. This library is explicitly designed for integration in Python pipelines, hiding the complexity of Tensor Core usage~\cite{Romein2021}.

Resource requirements in the chosen architecture are modest. Even a single modern desktop CPU can sample and FFT four data streams at several hundred Msps, and a mid-range GPU can perform all six cross-correlations in real time. For example, the CHIME Pathfinder X-engine used 16 AMD GPUs for 256 inputs~\cite{Denman2015}; by linear scaling, one GPU can handle well over four inputs. Callanan’s prototype likewise showed that one GPU (1080Ti) replaced two SKARAB boards (eight X-engines)~\cite{Callanan2020}. On the software side, DiFX running on ~1360 CPU cores handled a 10-antenna array~\cite{Deller2011}, so a trivial fraction of that power would serve \( N = 4 \). In summary, the correlator can run comfortably on a single server (8–16 CPU cores and 1–2 GPUs) without bottlenecks.

In conclusion, the literature strongly supports the design choices made. The analyses of correlator architectures (XF, FX, hybrid) and the case studies of ALMA, MeerKAT, DiFX, CHIME, and others provide a coherent picture: for a modest four-element array, an FX correlator is computationally appropriate, and software implementation is practical. Python, as the orchestrating language, brings flexibility and accessibility without insurmountable performance penalties, thanks to modern libraries and hardware. Resource requirements (CPU cores, GPUs, I/O bandwidth) will be modest, and latency demands can be met via buffering and parallelism. Therefore, a Python-based FX software correlator is fully justified for the intended system. It will offer the adaptability and integration ease required in a small observatory, while leveraging the high-level advantages of Python and the high-performance potential of GPUs.